\documentclass[a4paper,11pt]{article}

\usepackage[T1]{fontenc}
\usepackage[french]{babel}
\usepackage{geometry}
\usepackage{siunitx}
\sisetup{locale = FR, detect-all, inter-unit-product = \ensuremath{{}\cdot{}}}

\usepackage{tabavcmt}

\title{Package \texttt{tabavcmt}}
\author{Jean-Baptiste POUGET}
\date{Version 1.0 -- Juin 2026}

\begin{document}

\maketitle

\section{Présentation}

Le package \texttt{tabavcmt} permet de générer automatiquement des tableaux d'avancement en chimie à partir d'une syntaxe simple et intuitive.

Il est destiné à la rédaction de cours, d'exercices, de travaux dirigés et de corrigés de physique-chimie au lycée et dans l'enseignement supérieur.

Le package construit automatiquement :

\begin{itemize}
\item l'équation chimique ;
\item les séparations entre réactifs et produits ;
\item les lignes d'évolution des quantités de matière ;
\item les états du système ;
\item les expressions littérales en fonction de l'avancement ;
\item la mise en forme complète du tableau.
\end{itemize}

Les espèces chimiques sont saisies une seule fois sous la forme \verb|\ce{...}|. Le package extrait automatiquement le nom de l'espèce pour les quantités de matière en supprimant les états physiques \verb|(s)|, \verb|(l)|, \verb|(g)| et \verb|(aq)|.

\section{Chargement du package}

Le package se charge de manière classique : \begin{verbatim}
\usepackage{tabavcmt}
\end{verbatim}

\section{Environnement principal}

Les tableaux d'avancement sont créés à l'aide de l'environnement :

\begin{verbatim}
\begin{tabavcmt}[<options>]
Espèces chimiques
\end{tabavcmt}
\end{verbatim}

\subsection{Options disponibles}

\begin{center}
\begin{tabular}{|l|l|p{8cm}|}
\hline
Option & Valeur par défaut & Description \\
\hline
\texttt{fill} & \texttt{gray!10} & Couleur des en-têtes du tableau \\
\hline
\texttt{draw} & \texttt{black} & Couleur des bordures \\
\hline
\texttt{font size} & \verb|\normalsize| & Taille de la police \\
\hline
\texttt{scale} & \texttt{1} & Largeur maximale du tableau (adaptation automatique) \\
\hline
\texttt{xi} & \texttt{1} & Utilise $\xi$ comme variable d'avancement \\
\hline
\texttt{xi=0} & --- & Utilise $x$ comme variable d'avancement \\
\hline
\texttt{lignetotale} & \texttt{oui} & Affiche la ligne « État final (si totale) » \\
\hline
\texttt{lignetotale=non} & --- & Supprime cette ligne et remplace $\xi_f < \xi_{\max}$ par $\xi_f \leq \xi_{\max}$ \\
\hline
\texttt{linewidth} & 1pt & Ajuste l'épaisseur des lignes du tableau \\
\hline
\texttt{separation} & double & Sépare par une double ou simple ligne les réactifs des produits et la ligne optionnelle de l'état final dans le cas d'une réaction totale \\
\hline
\texttt{inter dbl lgn} & 1.2 & Gère la distance entre les lignes dans le cas d'une séparation par ligne double (clé \texttt{separation} vaut \texttt{double}) \\
\hline
\end{tabular}
\end{center}

\subsection{Déclaration des espèces chimiques}

\subsubsection{Réactif}

\begin{verbatim}
\reactif[<largeur>]{<coefficient>}{<formule chimique>}{<quantité initiale>}
\end{verbatim}

\subsubsection{Produit}

\begin{verbatim}
\produit[<largeur>]{<coefficient>}{<formule chimique>}{<quantité initiale>}
\end{verbatim}

L'argument \texttt{<largeur>} est optionnel. Sa valeur par défaut est :

\begin{verbatim}
[3cm]
\end{verbatim}

Lorsque le coefficient stœchiométrique vaut 1, il peut être omis :

\begin{verbatim}
\reactif{}{\ce{Fe2O3(s)}}{n_0}
\end{verbatim}

\section{Particularités}

\begin{itemize}
\item Les coefficients stœchiométriques sont automatiquement espacés des espèces chimiques.
\item Les états physiques sont conservés dans l'équation chimique.
\item Les états physiques sont automatiquement supprimés dans les expressions de type $n(\cdot)$.
\item La largeur du tableau est automatiquement ajustée grâce au package \texttt{adjustbox}.
\item Les petits tableaux ne sont jamais agrandis.
\item Les grands tableaux sont automatiquement réduits afin de tenir dans la largeur disponible.
\end{itemize}

\section{Exemples}
\subsection{Exemple 1}
\begin{verbatim}
\begin{tabavcmt}[font size = \normalsize,
fill=red,
draw=green,
xi=1,
scale=0.7,
ligne totale=non,
linewidth=1pt,
inter dbl lgn=1.2,
separation=simple
]
\reactif[2.5cm]{3}{\ce{Cl2(g)}}{\num{0,3}}
\reactif[2.5cm]{2}{\ce{Al(s)}}{\num{0,1}}
\produit[2.5cm]{2}{\ce{AlCl3(s)}}{0}
\end{tabavcmt}
\end{verbatim}

\begin{tabavcmt}[font size = \normalsize,
fill=red,
draw=green,
xi=0,
scale=0.9,
ligne totale=non,
linewidth=1pt,
inter dbl lgn=1.2,
separation=simple
]
\reactif[2.5cm]{3}{\ce{Cl2(g)}}{\num{0,3}}
\reactif[2.5cm]{2}{\ce{Al(s)}}{\num{0,1}}
\produit[2.5cm]{2}{\ce{AlCl3(s)}}{0}
\end{tabavcmt}

\newpage
\subsection{Exemple 2}
\begin{verbatim}
\begin{tabavcmt}[font size = \normalsize,
fill=blue!30,
draw=blue,
xi=0,
scale=1,
ligne totale=oui,
linewidth=0.5pt,
inter dbl lgn=1.5,
separation=double
]
\reactif[3.75cm]{}{\ce{Fe2O3(s)}}{\num{7,52e-2}}
\reactif[3.75cm]{2}{\ce{Al(s)}}{\num{1,48e-1}}
\produit[3cm]{2}{\ce{Fe(s)}}{0}
\produit[3cm]{}{\ce{Al2O3(s)}}{0}
\end{tabavcmt}
\end{verbatim}

\begin{tabavcmt}[font size = \normalsize,
fill=blue!30,
draw=blue,
xi=1,
scale=1,
ligne totale=oui,
linewidth=0.5pt,
inter dbl lgn=1.5,
separation=double
]
\reactif[3.75cm]{}{\ce{Fe2O3(s)}}{\num{7,52e-2}}
\reactif[3.75cm]{2}{\ce{Al(s)}}{\num{1,48e-1}}
\produit[3cm]{2}{\ce{Fe(s)}}{0}
\produit[3cm]{}{\ce{Al2O3(s)}}{0}
\end{tabavcmt}

\section{Compatibilité}

Le package a été développé et testé sous TeX Live 2019.

Les dépendances utilisées ont été volontairement limitées afin d'assurer une bonne compatibilité avec différentes versions de TeX Live.

\section{Historique des versions}

\begin{tabular}{ll}
1.0 & Création du package \\
\end{tabular}
\end{document}





